\section{早明的区域比较与制度实践}

同一项制度在不同地区从不以同一速度、同一成本或同一社会后果展开。账本中的诏令、军报、赋役记录、工程奏报与使节文书提供跨区域比较的起点；地方志、契约、碑刻、遗址与异域材料则提醒我们，中央分类无法覆盖地方实践。本节以事件链而非单一结论呈现制度、文化、环境与边地关系的差异。

\subsection{江南、华北与户役}

\paragraph{1398（洪武三十一年）五月24日：洪武帝生病了} 这一事件使区域差异进入可核查的编年。账本确认的范围是编年候选的年序可由官修与后出材料交叉；具体月日、参与者、战果或数字必须回查所列材料，不能将单一叙事当作定论。，但各材料服务于不同的行政、叙事或保存目的。事件之所以在该时点发生，与；可见的后续变化是它影响后续人事与政策议程；动机和派系标签不能由单一叙事裁定。比较不同地区时，不能把中心政策、地方报数、物质遗存与社会经验混为同一种证据。译写候选以编年材料定位；实录记载反映朝廷选择，崇祯期须另以奏疏、地方及敌对政权材料交叉。

\paragraph{1398（洪武三十一年）六月24日：洪武帝驾崩} 这一事件使区域差异进入可核查的编年。账本确认的范围是编年候选的年序可由官修与后出材料交叉；具体月日、参与者、战果或数字必须回查所列材料，不能将单一叙事当作定论。，但各材料服务于不同的行政、叙事或保存目的。事件之所以在该时点发生，与；可见的后续变化是它影响后续人事与政策议程；动机和派系标签不能由单一叙事裁定。比较不同地区时，不能把中心政策、地方报数、物质遗存与社会经验混为同一种证据。译写候选以编年材料定位；实录记载反映朝廷选择，崇祯期须另以奏疏、地方及敌对政权材料交叉。

\paragraph{1398（洪武三十一年）六月30日：朱允炆登基为建文帝} 这一事件使区域差异进入可核查的编年。账本确认的范围是编年候选的年序可由官修与后出材料交叉；具体月日、参与者、战果或数字必须回查所列材料，不能将单一叙事当作定论。，但各材料服务于不同的行政、叙事或保存目的。事件之所以在该时点发生，与；可见的后续变化是它影响后续人事与政策议程；动机和派系标签不能由单一叙事裁定。比较不同地区时，不能把中心政策、地方报数、物质遗存与社会经验混为同一种证据。译写候选以编年材料定位；实录记载反映朝廷选择，崇祯期须另以奏疏、地方及敌对政权材料交叉。

\paragraph{1398（洪武三十一年）：建文帝灭朱珪、朱勃、朱辅、朱翩诸侯国} 这一事件使区域差异进入可核查的编年。账本确认的范围是编年候选的年序可由官修与后出材料交叉；具体月日、参与者、战果或数字必须回查所列材料，不能将单一叙事当作定论。，但各材料服务于不同的行政、叙事或保存目的。事件之所以在该时点发生，与；可见的后续变化是它影响后续人事与政策议程；动机和派系标签不能由单一叙事裁定。比较不同地区时，不能把中心政策、地方报数、物质遗存与社会经验混为同一种证据。译写候选以编年材料定位；实录记载反映朝廷选择，崇祯期须另以奏疏、地方及敌对政权材料交叉。

\paragraph{1398（洪武三十一年）：中国最后一次活人祭祀记录} 这一事件使区域差异进入可核查的编年。账本确认的范围是编年候选的年序可由官修与后出材料交叉；具体月日、参与者、战果或数字必须回查所列材料，不能将单一叙事当作定论。，但各材料服务于不同的行政、叙事或保存目的。事件之所以在该时点发生，与；可见的后续变化是它影响后续人事与政策议程；动机和派系标签不能由单一叙事裁定。比较不同地区时，不能把中心政策、地方报数、物质遗存与社会经验混为同一种证据。译写候选以编年材料定位；实录记载反映朝廷选择，崇祯期须另以奏疏、地方及敌对政权材料交叉。

\paragraph{1399（建文元年）六月：建文帝归还朱棣的儿子} 这一事件使区域差异进入可核查的编年。账本确认的范围是编年候选的年序可由官修与后出材料交叉；具体月日、参与者、战果或数字必须回查所列材料，不能将单一叙事当作定论。，但各材料服务于不同的行政、叙事或保存目的。事件之所以在该时点发生，与；可见的后续变化是它影响后续人事与政策议程；动机和派系标签不能由单一叙事裁定。比较不同地区时，不能把中心政策、地方报数、物质遗存与社会经验混为同一种证据。译写候选以编年材料定位；实录记载反映朝廷选择，崇祯期须另以奏疏、地方及敌对政权材料交叉。

\subsection{西南、东北与边地中介}

\paragraph{1399（建文元年）七月：一名军官以煽动叛乱罪逮捕了朱棣的两名下级官员} 这一事件使区域差异进入可核查的编年。账本确认的范围是编年候选的年序可由官修与后出材料交叉；具体月日、参与者、战果或数字必须回查所列材料，不能将单一叙事当作定论。，但各材料服务于不同的行政、叙事或保存目的。事件之所以在该时点发生，与；可见的后续变化是它改变了后续调兵、补给或地方秩序的条件；战果和伤亡须分开核对。比较不同地区时，不能把中心政策、地方报数、物质遗存与社会经验混为同一种证据。译写候选以编年材料定位；实录记载反映朝廷选择，崇祯期须另以奏疏、地方及敌对政权材料交叉。

\paragraph{1399（建文元年）八月5日：荆南战役：朱棣向邻县发起进攻} 这一事件使区域差异进入可核查的编年。账本确认的范围是编年候选的年序可由官修与后出材料交叉；具体月日、参与者、战果或数字必须回查所列材料，不能将单一叙事当作定论。，但各材料服务于不同的行政、叙事或保存目的。事件之所以在该时点发生，与；可见的后续变化是它改变了后续调兵、补给或地方秩序的条件；战果和伤亡须分开核对。比较不同地区时，不能把中心政策、地方报数、物质遗存与社会经验混为同一种证据。译写候选以编年材料定位；实录记载反映朝廷选择，崇祯期须另以奏疏、地方及敌对政权材料交叉。

\paragraph{1399（建文元年）九月25日：靖难之役：朱棣大败建文帝十三万大军} 这一事件使区域差异进入可核查的编年。账本确认的范围是编年候选的年序可由官修与后出材料交叉；具体月日、参与者、战果或数字必须回查所列材料，不能将单一叙事当作定论。，但各材料服务于不同的行政、叙事或保存目的。事件之所以在该时点发生，与；可见的后续变化是它改变了后续调兵、补给或地方秩序的条件；战果和伤亡须分开核对。比较不同地区时，不能把中心政策、地方报数、物质遗存与社会经验混为同一种证据。译写候选以编年材料定位；实录记载反映朝廷选择，崇祯期须另以奏疏、地方及敌对政权材料交叉。

\paragraph{1399（建文元年）十一月12日：靖难之役：建文帝军队围攻北平，三周后被迫撤退} 这一事件使区域差异进入可核查的编年。账本确认的范围是编年候选的年序可由官修与后出材料交叉；具体月日、参与者、战果或数字必须回查所列材料，不能将单一叙事当作定论。，但各材料服务于不同的行政、叙事或保存目的。事件之所以在该时点发生，与；可见的后续变化是它改变了后续调兵、补给或地方秩序的条件；战果和伤亡须分开核对。比较不同地区时，不能把中心政策、地方报数、物质遗存与社会经验混为同一种证据。译写候选以编年材料定位；实录记载反映朝廷选择，崇祯期须另以奏疏、地方及敌对政权材料交叉。

\paragraph{1400（建文二年）一月：靖难之役：朱棣攻入山西} 这一事件使区域差异进入可核查的编年。账本确认的范围是编年候选的年序可由官修与后出材料交叉；具体月日、参与者、战果或数字必须回查所列材料，不能将单一叙事当作定论。，但各材料服务于不同的行政、叙事或保存目的。事件之所以在该时点发生，与；可见的后续变化是它改变了后续调兵、补给或地方秩序的条件；战果和伤亡须分开核对。比较不同地区时，不能把中心政策、地方报数、物质遗存与社会经验混为同一种证据。译写候选以编年材料定位；实录记载反映朝廷选择，崇祯期须另以奏疏、地方及敌对政权材料交叉。

\paragraph{1400（建文二年）五月18日：靖难之战：朱棣军伤亡惨重} 这一事件使区域差异进入可核查的编年。账本确认的范围是编年候选的年序可由官修与后出材料交叉；具体月日、参与者、战果或数字必须回查所列材料，不能将单一叙事当作定论。，但各材料服务于不同的行政、叙事或保存目的。事件之所以在该时点发生，与；可见的后续变化是它改变了后续调兵、补给或地方秩序的条件；战果和伤亡须分开核对。比较不同地区时，不能把中心政策、地方报数、物质遗存与社会经验混为同一种证据。译写候选以编年材料定位；实录记载反映朝廷选择，崇祯期须另以奏疏、地方及敌对政权材料交叉。

\subsection{城市、道路与文化资源}

\paragraph{1400（建文二年）六月8日：靖难之役：朱棣围攻德州} 这一事件使区域差异进入可核查的编年。账本确认的范围是编年候选的年序可由官修与后出材料交叉；具体月日、参与者、战果或数字必须回查所列材料，不能将单一叙事当作定论。，但各材料服务于不同的行政、叙事或保存目的。事件之所以在该时点发生，与；可见的后续变化是它改变了后续调兵、补给或地方秩序的条件；战果和伤亡须分开核对。比较不同地区时，不能把中心政策、地方报数、物质遗存与社会经验混为同一种证据。译写候选以编年材料定位；实录记载反映朝廷选择，崇祯期须另以奏疏、地方及敌对政权材料交叉。

\paragraph{1400（建文二年）九月4日：靖难之役：朱棣解除德州之围返回北平} 这一事件使区域差异进入可核查的编年。账本确认的范围是编年候选的年序可由官修与后出材料交叉；具体月日、参与者、战果或数字必须回查所列材料，不能将单一叙事当作定论。，但各材料服务于不同的行政、叙事或保存目的。事件之所以在该时点发生，与；可见的后续变化是它改变了后续调兵、补给或地方秩序的条件；战果和伤亡须分开核对。比较不同地区时，不能把中心政策、地方报数、物质遗存与社会经验混为同一种证据。译写候选以编年材料定位；实录记载反映朝廷选择，崇祯期须另以奏疏、地方及敌对政权材料交叉。

\paragraph{1401（建文三年）一月9日：靖难之役：朱棣部队在山东遭炸，伤亡惨重，被迫撤退} 这一事件使区域差异进入可核查的编年。账本确认的范围是编年候选的年序可由官修与后出材料交叉；具体月日、参与者、战果或数字必须回查所列材料，不能将单一叙事当作定论。，但各材料服务于不同的行政、叙事或保存目的。事件之所以在该时点发生，与；可见的后续变化是它改变了后续调兵、补给或地方秩序的条件；战果和伤亡须分开核对。比较不同地区时，不能把中心政策、地方报数、物质遗存与社会经验混为同一种证据。译写候选以编年材料定位；实录记载反映朝廷选择，崇祯期须另以奏疏、地方及敌对政权材料交叉。

\paragraph{1401（建文三年）四月5日：靖难之役：朱棣军队在德州附近大败皇军} 这一事件使区域差异进入可核查的编年。账本确认的范围是编年候选的年序可由官修与后出材料交叉；具体月日、参与者、战果或数字必须回查所列材料，不能将单一叙事当作定论。，但各材料服务于不同的行政、叙事或保存目的。事件之所以在该时点发生，与；可见的后续变化是它改变了后续调兵、补给或地方秩序的条件；战果和伤亡须分开核对。比较不同地区时，不能把中心政策、地方报数、物质遗存与社会经验混为同一种证据。译写候选以编年材料定位；实录记载反映朝廷选择，崇祯期须另以奏疏、地方及敌对政权材料交叉。

\paragraph{1401（建文三年）八月：靖难之战：皇军迫使朱棣北撤北平} 这一事件使区域差异进入可核查的编年。账本确认的范围是编年候选的年序可由官修与后出材料交叉；具体月日、参与者、战果或数字必须回查所列材料，不能将单一叙事当作定论。，但各材料服务于不同的行政、叙事或保存目的。事件之所以在该时点发生，与；可见的后续变化是它改变了后续调兵、补给或地方秩序的条件；战果和伤亡须分开核对。比较不同地区时，不能把中心政策、地方报数、物质遗存与社会经验混为同一种证据。译写候选以编年材料定位；实录记载反映朝廷选择，崇祯期须另以奏疏、地方及敌对政权材料交叉。

\paragraph{1401（建文三年）八月：建文帝限制佛教和道教的土地占有规模} 这一事件使区域差异进入可核查的编年。账本确认的范围是编年候选的年序可由官修与后出材料交叉；具体月日、参与者、战果或数字必须回查所列材料，不能将单一叙事当作定论。，但各材料服务于不同的行政、叙事或保存目的。事件之所以在该时点发生，与；可见的后续变化是它影响后续人事与政策议程；动机和派系标签不能由单一叙事裁定。比较不同地区时，不能把中心政策、地方报数、物质遗存与社会经验混为同一种证据。译写候选以编年材料定位；实录记载反映朝廷选择，崇祯期须另以奏疏、地方及敌对政权材料交叉。

\subsection{环境、工程与地方应对}

\paragraph{1401（建文三年）十月：靖难之战：帝国军队被驱逐出北平地区} 这一事件使区域差异进入可核查的编年。账本确认的范围是编年候选的年序可由官修与后出材料交叉；具体月日、参与者、战果或数字必须回查所列材料，不能将单一叙事当作定论。，但各材料服务于不同的行政、叙事或保存目的。事件之所以在该时点发生，与；可见的后续变化是它改变了后续调兵、补给或地方秩序的条件；战果和伤亡须分开核对。比较不同地区时，不能把中心政策、地方报数、物质遗存与社会经验混为同一种证据。译写候选以编年材料定位；实录记载反映朝廷选择，崇祯期须另以奏疏、地方及敌对政权材料交叉。

\paragraph{1402（建文四年）一月：靖难之战：朱棣攻克山东西北} 这一事件使区域差异进入可核查的编年。账本确认的范围是编年候选的年序可由官修与后出材料交叉；具体月日、参与者、战果或数字必须回查所列材料，不能将单一叙事当作定论。，但各材料服务于不同的行政、叙事或保存目的。事件之所以在该时点发生，与；可见的后续变化是它改变了后续调兵、补给或地方秩序的条件；战果和伤亡须分开核对。比较不同地区时，不能把中心政策、地方报数、物质遗存与社会经验混为同一种证据。译写候选以编年材料定位；实录记载反映朝廷选择，崇祯期须另以奏疏、地方及敌对政权材料交叉。

\paragraph{1402（建文四年）三月3日：靖难之役：朱棣攻克徐州} 这一事件使区域差异进入可核查的编年。账本确认的范围是编年候选的年序可由官修与后出材料交叉；具体月日、参与者、战果或数字必须回查所列材料，不能将单一叙事当作定论。，但各材料服务于不同的行政、叙事或保存目的。事件之所以在该时点发生，与；可见的后续变化是它改变了后续调兵、补给或地方秩序的条件；战果和伤亡须分开核对。比较不同地区时，不能把中心政策、地方报数、物质遗存与社会经验混为同一种证据。译写候选以编年材料定位；实录记载反映朝廷选择，崇祯期须另以奏疏、地方及敌对政权材料交叉。

\paragraph{1402（建文四年）四月：靖难之役：朱棣在苏州大败皇军} 这一事件使区域差异进入可核查的编年。账本确认的范围是编年候选的年序可由官修与后出材料交叉；具体月日、参与者、战果或数字必须回查所列材料，不能将单一叙事当作定论。，但各材料服务于不同的行政、叙事或保存目的。事件之所以在该时点发生，与；可见的后续变化是它改变了后续调兵、补给或地方秩序的条件；战果和伤亡须分开核对。比较不同地区时，不能把中心政策、地方报数、物质遗存与社会经验混为同一种证据。译写候选以编年材料定位；实录记载反映朝廷选择，崇祯期须另以奏疏、地方及敌对政权材料交叉。

\paragraph{1402（建文四年）五月23日：靖难之役：朱棣在安徽被皇军击退} 这一事件使区域差异进入可核查的编年。账本确认的范围是编年候选的年序可由官修与后出材料交叉；具体月日、参与者、战果或数字必须回查所列材料，不能将单一叙事当作定论。，但各材料服务于不同的行政、叙事或保存目的。事件之所以在该时点发生，与；可见的后续变化是它改变了后续调兵、补给或地方秩序的条件；战果和伤亡须分开核对。比较不同地区时，不能把中心政策、地方报数、物质遗存与社会经验混为同一种证据。译写候选以编年材料定位；实录记载反映朝廷选择，崇祯期须另以奏疏、地方及敌对政权材料交叉。

\paragraph{1402（建文四年）五月28日：靖难之役：朱棣在灵璧大败皇军} 这一事件使区域差异进入可核查的编年。账本确认的范围是编年候选的年序可由官修与后出材料交叉；具体月日、参与者、战果或数字必须回查所列材料，不能将单一叙事当作定论。，但各材料服务于不同的行政、叙事或保存目的。事件之所以在该时点发生，与；可见的后续变化是它改变了后续调兵、补给或地方秩序的条件；战果和伤亡须分开核对。比较不同地区时，不能把中心政策、地方报数、物质遗存与社会经验混为同一种证据。译写候选以编年材料定位；实录记载反映朝廷选择，崇祯期须另以奏疏、地方及敌对政权材料交叉。

\subsection{环境、工程与地方应对}

\paragraph{1402（建文四年）六月7日：靖难之役：朱棣大军渡过淮河} 这一事件使区域差异进入可核查的编年。账本确认的范围是编年候选的年序可由官修与后出材料交叉；具体月日、参与者、战果或数字必须回查所列材料，不能将单一叙事当作定论。，但各材料服务于不同的行政、叙事或保存目的。事件之所以在该时点发生，与；可见的后续变化是其法定安排不等于地方执行，后续须以地域材料追踪负担与成效。比较不同地区时，不能把中心政策、地方报数、物质遗存与社会经验混为同一种证据。译写候选以编年材料定位；实录记载反映朝廷选择，崇祯期须另以奏疏、地方及敌对政权材料交叉。

\paragraph{1402（建文四年）六月17日：靖难之役：朱棣拿下扬州} 这一事件使区域差异进入可核查的编年。账本确认的范围是编年候选的年序可由官修与后出材料交叉；具体月日、参与者、战果或数字必须回查所列材料，不能将单一叙事当作定论。，但各材料服务于不同的行政、叙事或保存目的。事件之所以在该时点发生，与；可见的后续变化是它改变了后续调兵、补给或地方秩序的条件；战果和伤亡须分开核对。比较不同地区时，不能把中心政策、地方报数、物质遗存与社会经验混为同一种证据。译写候选以编年材料定位；实录记载反映朝廷选择，崇祯期须另以奏疏、地方及敌对政权材料交叉。

\paragraph{1402（建文四年）七月1日：靖难之役：朱棣被拦在南京对面的长江边} 这一事件使区域差异进入可核查的编年。账本确认的范围是编年候选的年序可由官修与后出材料交叉；具体月日、参与者、战果或数字必须回查所列材料，不能将单一叙事当作定论。，但各材料服务于不同的行政、叙事或保存目的。事件之所以在该时点发生，与；可见的后续变化是它改变了后续调兵、补给或地方秩序的条件；战果和伤亡须分开核对。比较不同地区时，不能把中心政策、地方报数、物质遗存与社会经验混为同一种证据。译写候选以编年材料定位；实录记载反映朝廷选择，崇祯期须另以奏疏、地方及敌对政权材料交叉。

\paragraph{1402（建文四年）七月3日：靖难之役：布政司侍郎陈选投奔朱棣，起义军渡过长江} 这一事件使区域差异进入可核查的编年。账本确认的范围是编年候选的年序可由官修与后出材料交叉；具体月日、参与者、战果或数字必须回查所列材料，不能将单一叙事当作定论。，但各材料服务于不同的行政、叙事或保存目的。事件之所以在该时点发生，与；可见的后续变化是它改变了后续调兵、补给或地方秩序的条件；战果和伤亡须分开核对。比较不同地区时，不能把中心政策、地方报数、物质遗存与社会经验混为同一种证据。译写候选以编年材料定位；实录记载反映朝廷选择，崇祯期须另以奏疏、地方及敌对政权材料交叉。

\paragraph{1402（建文四年）七月13日：靖难之战：朱晖打开南京金川门，不战而让朱棣入城；建文帝失踪，家人被监禁} 这一事件使区域差异进入可核查的编年。账本确认的范围是编年候选的年序可由官修与后出材料交叉；具体月日、参与者、战果或数字必须回查所列材料，不能将单一叙事当作定论。，但各材料服务于不同的行政、叙事或保存目的。事件之所以在该时点发生，与；可见的后续变化是它改变了后续调兵、补给或地方秩序的条件；战果和伤亡须分开核对。比较不同地区时，不能把中心政策、地方报数、物质遗存与社会经验混为同一种证据。译写候选以编年材料定位；实录记载反映朝廷选择，崇祯期须另以奏疏、地方及敌对政权材料交叉。

\paragraph{1402（建文四年）七月17日：朱棣即位，是为永乐皇帝} 这一事件使区域差异进入可核查的编年。账本确认的范围是编年候选的年序可由官修与后出材料交叉；具体月日、参与者、战果或数字必须回查所列材料，不能将单一叙事当作定论。，但各材料服务于不同的行政、叙事或保存目的。事件之所以在该时点发生，与；可见的后续变化是它影响后续人事与政策议程；动机和派系标签不能由单一叙事裁定。比较不同地区时，不能把中心政策、地方报数、物质遗存与社会经验混为同一种证据。译写候选以编年材料定位；实录记载反映朝廷选择，崇祯期须另以奏疏、地方及敌对政权材料交叉。
